\section{OpenAIの生命科学モデル名が価格ページに登場}
OpenAIのAPI価格ページで、生命科学向けの「gpt-rosalind-discovery」が既存の「gpt-rosalind-research」と並んで記載されていると話題になった。Grokは両モデルの入力・出力単価が100万トークン当たり5ドル・25ドルであることを確認したと記録する。一方、公式ブログやシステムカードは確認されておらず、提供条件や両モデルの機能差は不明。新モデルの正式発表と同一視できない。
\dailyxsource{AI Now (@AI\_N0W)}{https://x.com/AI_N0W/status/2109393758920032384}

\section{EmbeddingGemma 2のGGUF版とオンデバイス利用}
Googleが10月6日に公開したEmbeddingGemma 2について、GGUF量子化版の共有やローカル実行の話題が続いた。740Mパラメータのマルチモーダル埋め込みモデルで、テキスト・画像・音声・動画を共通の埋め込み空間に扱うとされる。プライバシーや運用費用の面でオンデバイス処理への期待があるが、リリース自体は観測窓外。利用者の性能報告を独立したベンチマーク結果とは扱わない。
\dailyxsource{Hugging Models (@HuggingModels)}{https://x.com/HuggingModels/status/2109393610693321212}

\section{Claude MotionとDashboardsのベータが話題}
ClaudeのMotionとDashboardsが継続的に紹介された。Motionはコードに基づいて編集可能なアニメーションを生成し、MP4へ書き出す機能とされ、動画生成モデルとは異なる仕組みが特徴。Dashboardsはデータプラットフォーム連携を含む企業向け機能として言及された。公式リソースは10月8日頃の観測窓外で、提供対象はTeam・Enterprise向けとされる。品質や正式提供時期は未確認。
\dailyxsource{The AI Daily Brief (@AIDailyBrief)}{https://x.com/AIDailyBrief/status/2109388617869328554}

\section{RSIArenaの第1段階結果に関する投稿}
再帰的自己改善を競うRSIArenaのStage 1が終了し、Grok 4.7が訓練したモデルが上位に入ったとする投稿があった。投稿者は、10モデルによる競争、1,685票のブラインド比較、APIクレジット不足による早期終了などを挙げている。公式の順位表や評価条件は確認されていない。\textbf{根拠投稿は07:04 JSTで、今回の観測窓終了後に当たる。} 本号では原文収録の参考情報として扱う。
\dailyxsource{Zichen Chen (@my\_cat\_can\_code)}{https://x.com/my_cat_can_code/status/2109404695907410162}

\section{Qwen画像モデルTurbo版の高速化に言及}
Qwenの画像生成モデルに、生成ステップ数を従来の40から8へ減らしたTurbo版があるとの投稿が紹介された。7B規模、2K解像度、研究用ライセンスなどが主張され、別投稿ではNVFP4量子化版のローカル実行も話題になった。ただし公式リポジトリ、ライセンス、画質や速度の比較条件は未確認。\textbf{主要投稿は07:02 JSTで観測窓外}のため、確定した当日発表としては扱わない。
\dailyxsource{Context Window (@ContextWindow)}{https://x.com/ContextWindow/status/2109404314812965083}

\section{WebMCPの標準化議論とMCP Gateway}
Webサイトの操作機能をエージェントにスキーマで公開するWebMCPの仕様議論と、ツール群を整理する自己ホスト型MCP Gatewayが紹介された。WebMCPはブラウザーの実験的な実装に言及があり、GatewayはRust製でツールのオンデマンド発見やOpenAPI取り込みを掲げる。後者のライセンスはPolyForm Noncommercialとされる。どちらも実装・普及の段階を確認したものではなく、投稿に基づく技術動向の紹介である。
\dailyxsource{Prasenjit Sarkar (@stretchcloud)}{https://x.com/stretchcloud/status/2109293443268948142}

\section{Sakana AIの多層論文レビュー研究}
Sakana AIの「Multi-Layered Review」が、複数のClaudeエージェントを使って論文の問題点を見つける仕組みとして紹介された。投稿では、PDFの図や数式を含めて処理し、ベンチマークで中心的主張の矛盾の73\%を検出、1回約47セントと説明されている。論文本文や評価手法はGrok収集で確認できておらず、数値は紹介投稿の主張にとどまる。査読支援の自動化という方向性が注目点となる。
\dailyxsource{AIQUEST (@AiquestAcademy)}{https://x.com/AiquestAcademy/status/2109148632351916511}

\section{長時間タスクにおけるエージェントハーネス評価}
17名の研究者によるとされる論文が、エージェントの長時間タスクでの失敗をモデル単体の能力だけでなく、観測・制御ループや安全機構などのハーネスとの相互作用から分析した、と紹介された。モデル、ランタイム、タスク構造、評価設計を一体で考える必要性が示唆される。ただし論文URLや実験詳細は未確認で、特定の失敗原因が一般に支配的だと結論付けることはできない。
\dailyxsource{slash1s (@slash1sol)}{https://x.com/slash1sol/status/2109383952850993419}

\section{Amazonの研究ワールドモデル論文を紹介}
Amazonの研究として、過去の実験ログを参照するLLMが、学習条件の変更による結果を予測する「研究ワールドモデル」の取り組みが共有された。紹介投稿は2,653件の実験記録を用い、平均ランキング相関が0.506から0.774へ改善したとする。成功例だけでなく失敗実験の記録も活用する点が関心を集めたが、論文本文や再現条件は確認されていない。数字は投稿者の説明として記録する。
\dailyxsource{Rohan Paul (@rohanpaul\_ai)}{https://x.com/rohanpaul_ai/status/2109243398628708619}

\section{物理AI・ロボティクス関連ニュースのまとめ}
Boston DynamicsのCEO人事、ロボティクス企業の資金調達、Odyssey-3ワールドモデル、NVIDIAのVideo to Data Challengeなどをまとめた投稿が収集された。実世界データやロボットの運用基盤への関心を示す一方、各社の公式発表は今回確認されていない。\textbf{まとめ投稿の時刻は07:02 JSTで観測窓終了後}に当たるため、個々の出来事が当日発表されたとの判断は避ける。
\dailyxsource{Continuum Labs (@continuumlabs\_)}{https://x.com/continuumlabs_/status/2109404380751704478}

\section{Microsoft-Decision-1の公開をめぐる言及}
MicrosoftがQwen3.5-9Bを追加学習した意思決定モデル「Microsoft-Decision-1」を公開したとのダイジェスト投稿があった。構造化された選択肢へ確率を返す用途や、Foundry・OpenRouter経由の提供、36種類のベンチマークでの性能が主張されている。ただし公式ドキュメントや比較条件は確認されておらず、低レイテンシなどの優位性も投稿内容の紹介にとどめる。
\dailyxsource{Mike T. (@daogangtang)}{https://x.com/daogangtang/status/2109348572491841753}

\section{Falcon ASRという音声認識モデルへの言及}
Hugging Faceが「Falcon ASR」という音声認識モデルを公開したとの短い投稿が収集された。ただしニュースまとめの一部として言及されたのみで、モデルカード、ライセンス、性能や公開主体を裏付ける資料は未確認。\textbf{投稿時刻は07:02 JSTで観測窓外}である。現時点では新モデルの公開を確認したニュースではなく、今後の追跡候補として記録する。
\dailyxsource{Peter (@Peterjordenson)}{https://x.com/Peterjordenson/status/2109404229387354272}
